\section{Construction of the Rejector Loss}
\label{sec:method}

The Bayes rejector of Lemma~\ref{lem:bayes-rej} compares the conditional
risk $\eta_h(x)=\E[\psi(h(x),Y)\mid X=x]$ with the abstention cost $c(x)$.
For a PINN deployed on a deterministic boundary-value problem two further
steps are needed before this rule can be implemented. First, the abstract
discrepancy $\psi$ has to be instantiated in a way that aligns the
abstention layer with the pointwise notion of error practitioners report
(Section~\ref{sec:method-psi}). Second, $\eta_h$ is unobservable -- it
depends on the unknown solution $u^\star$, and obtaining labels would
require running the very fallback solver the PINN is meant to amortise --
so the rejector must be trained against a self-supervised proxy built
from quantities the PINN itself exposes
(Sections~\ref{sec:method-proxy}--\ref{sec:method-consistency}).

\subsection{The pointwise discrepancy in the PDE setting}
\label{sec:method-psi}

For a deterministic PDE the conditional law of $Y$ given $X=x$ is the
point mass at $u^\star(x)$, so $\eta_h(x)=\psi(h(x),u^\star(x))$ collapses
to a deterministic pointwise quantity. The choice of $\psi$ then directly
fixes which notion of error the rejector enforces. We take
\begin{equation}
  \psi(a,b)\;:=\;\|a-b\|_2,
  \label{eq:psi-norm}
\end{equation}
the unsquared Euclidean discrepancy. Two reasons recommend
\eqref{eq:psi-norm} over alternatives such as squared, relative, or
operator-weighted norms. First, $\eta_h(x)=\|h(x)-u^\star(x)\|_2$ is
exactly the pointwise solution error: a region where it is large is a
region whose PINN values would visibly contaminate downstream physical
diagnostics (forces, fluxes, integrated quantities), while a region where
it is small is one where the PINN can be deployed as drop-in for the
solver. The threshold rule of Lemma~\ref{lem:bayes-rej} therefore reads as
``replace the PINN exactly where its pointwise error exceeds the per-point
cost of running the solver instead'', which is the operational statement
the abstention layer is meant to provide. Second, with this choice $\eta_h$
and the proxy constructed below share the same dimensions and the same
units as $u^\star$, so the cost $c(x)$ is expressed in the same unit
system, and the threshold rule is dimensionally homogeneous across PDE
families and grids.

The substantive obstacle is that $u^\star$, and hence $\eta_h$, is
unavailable. The rejector therefore cannot be trained against
\eqref{eq:psi-norm} directly; it must be trained against a computable
proxy $\rho^\theta(x)$ that approximates $\eta_h(x)$ from quantities
derivable from $h$ alone. Constructing such a proxy is the subject of
Section~\ref{sec:method-proxy}; the rest of the method then proceeds with
$\rho^\theta$ playing the role $\eta_h$ would have played if it had been
observable.

\subsection{A self-supervised proxy from residual and error transport}
\label{sec:method-proxy}

The strong-form interior residual of the trained PINN $h=\widehat{u}_\theta$,
\begin{equation*}
  \mathfrak{r}^\theta(x)\;:=\;\Ncal[h](x)-f(x),\qquad x\in\Omega,
\end{equation*}
is the most natural pointwise imbalance signal and is computed by
automatic differentiation through $h$ at negligible cost. Used in
isolation, however, it is a poor proxy for $\eta_h$: the PDE operator
generally smooths or attenuates errors, allowing $|h(x)-u^\star(x)|$ to be
large where $\|\mathfrak{r}^\theta(x)\|$ is small and vice versa
\citep{verfurth2013posteriori,phillips2011residual}. This is the
well-known \emph{residual--error gap}, and it is precisely the gap that an
a-posteriori bound such as \eqref{eq:apost} closes only in a global norm.

The error-transport equation closes the gap pointwise by transporting the
residual through a linearisation of $\Ncal$ at $h$. Let $\Lcal_h$ denote
the formal Fr\'echet derivative of $\Ncal$ at $h$ (e.g.\ for steady
advection--diffusion $\Ncal[u]=u\!\cdot\!\nabla u-\nu\Delta u$,
$\Lcal_h[v]=h\!\cdot\!\nabla v+v\!\cdot\!\nabla h-\nu\Delta v$;
for the Poisson operator $\Lcal_h=-\Delta$ regardless of $h$). Define
the \emph{error-transport estimate} $e^\theta:\Omega\to\R^{d_u}$ as the
solution of the linearised boundary-value problem
\begin{equation}
  \Lcal_h[e^\theta](x)\;=\;\mathfrak{r}^\theta(x),
  \quad x\in\Omega,
  \qquad
  e^\theta\big|_{\partial\Omega}\;=\;0,
  \label{eq:ete}
\end{equation}
discretised on the same grid the fallback solver will use and computed by
a single Jacobi sweep, FFT inverse of $-\Delta$, or whatever direct
linear solve is most natural for $\Lcal_h$ -- in every case a cost
strictly dominated by one full nonlinear fallback solve. The Taylor
expansion $\Ncal[h]-\Ncal[u^\star]
=\Lcal_h[h-u^\star]+\Rcal_2(h-u^\star)$ together with $\Ncal[u^\star]=f$
gives $\Lcal_h[h-u^\star]=\mathfrak{r}^\theta+\Rcal_2(h-u^\star)$; combining
with~\eqref{eq:ete} yields
$e^\theta-(h-u^\star)=\Lcal_h^{-1}[-\Rcal_2(h-u^\star)]$, so $e^\theta$
recovers the pointwise error to first order in $\|h-u^\star\|$
\citep{phillips2011residual}.

\paragraph{Per-summand normalisation.}
The two natural pointwise signals $\|\mathfrak{r}^\theta(x)\|$ and
$\|e^\theta(x)\|$ are not directly commensurable. By construction
$e^\theta$ shares the units of $u^\star$, whereas $\mathfrak{r}^\theta=
\Ncal[h]-f$ inherits the units of $\Ncal[u^\star]$ and so differs from
$e^\theta$ by the units of the differential operator (e.g.\
$\mathrm{length}^{-2}$ for $-\Delta$). Adding them raw would couple the
threshold $c(x)$ to the choice of physical units and to the grid
spacing, which we want to avoid: a router calibrated on one PDE family or
discretisation should not have to be retuned every time units or grid
resolution change. We therefore rescale each summand by its own domain
median before combining. Let
\begin{equation}
  m_{\mathrm{r}}\;:=\;\mathrm{median}_{x\in\Omega}\,\bigl\|\mathfrak{r}^\theta(x)\bigr\|,
  \qquad
  m_{\mathrm{e}}\;:=\;\mathrm{median}_{x\in\Omega}\,\bigl\|e^\theta(x)\bigr\|,
  \label{eq:medians}
\end{equation}
both of which are strictly positive whenever $h\neq u^\star$ on a set of
positive measure (else the routing problem is empty). The proxy is then
\begin{equation}
  \rho^\theta(x)\;:=\;
  \frac{\bigl\|\mathfrak{r}^\theta(x)\bigr\|}{m_{\mathrm{r}}}
  \;+\;
  \frac{\bigl\|e^\theta(x)\bigr\|}{m_{\mathrm{e}}},
  \label{eq:rho-def}
\end{equation}
which is a dimensionless non-negative scalar field with
$\mathrm{median}_{x\in\Omega}\,\rho^\theta(x)\in[1,2]$ by construction
(each summand has unit median; their sum has median in this range by
sub-additivity of the median functional). Median rather than mean
normalisation is deliberate: PINN residual fields are heavy-tailed and
right-skewed \citep{wang2021understanding,wang2022whenpinns}, so the mean
is dominated by the same hot spots the rejector is meant to detect,
collapsing $\rho^\theta\approx 1$ everywhere; the median tracks the bulk
and pushes hot spots into the long tail $\rho^\theta\gg 1$ where the
threshold rule must operate. Throughout the rest of the paper, $c(x)$ is
expressed in the same dimensionless units as \eqref{eq:rho-def}, so the
threshold rule $\rho^\theta(x)>c(x)$ is invariant to unit and grid
changes that affect $\|\mathfrak{r}^\theta\|$ and $\|e^\theta\|$
multiplicatively.

\begin{assumption}[Linearisation faithfulness of the proxy]
\label{ass:linearisation}
There exist constants $\kappa_1,\kappa_2>0$ depending on the PDE and on
$\theta$, and a non-negative linearisation residual
$\delta_{\mathrm{lin}}^\theta:\Omega\to\R_{\geq 0}$ with
$\delta_{\mathrm{lin}}^\theta(x)=o(\|h(x)-u^\star(x)\|)$ as
$\|h-u^\star\|_{L^\infty}\to 0$, such that for $\Pr_X$-a.e.\ $x\in\Omega$,
\[
  \kappa_1\,\frac{\bigl\|h(x)-u^\star(x)\bigr\|}{m_{u}}
  \;\leq\;\rho^\theta(x)\;\leq\;
  \kappa_2\,\frac{\bigl\|h(x)-u^\star(x)\bigr\|}{m_{u}}
  \;+\;\delta_{\mathrm{lin}}^\theta(x),
\]
where $m_{u}:=\mathrm{median}_{x\in\Omega}\|h(x)-u^\star(x)\|$ is the
median pointwise PINN error.
\end{assumption}

The two summands of \eqref{eq:rho-def} carry complementary information.
$\|\mathfrak{r}^\theta\|/m_{\mathrm{r}}$ flags equation imbalance and is
local; it contributes the upper bound through the bounded-operator
estimate
$\|\mathfrak{r}^\theta\|\leq\|\Lcal_h\|\,\|h-u^\star\|+\|\Rcal_2\|$.
$\|e^\theta\|/m_{\mathrm{e}}$ contributes the lower bound: by~\eqref{eq:ete}
a non-zero residual at one location reaches far-field points
proportionally to the Green's function of $\Lcal_h$, so $\|e^\theta\|$
tracks the \emph{spread} of the imbalance through the operator and
prevents the proxy from collapsing to zero at points where the residual
happens to be small but the propagated error is not. The normalised sum
therefore flags both local imbalance and operator-mediated propagation,
in shared dimensionless units.

Substituting \eqref{eq:rho-def} for $\eta_h$ in the abstention loss
\eqref{eq:abs-loss} gives the \emph{computable} pointwise abstention loss
\begin{equation}
  L(s)
  \;:=\;
  \E_X\!\Bigl[\rho^\theta(X)\,\ind\{s(X)\leq 0\}+c(X)\,\ind\{s(X)>0\}\Bigr],
  \label{eq:proxy-target}
\end{equation}
parameterised by a measurable score $s\in\Hcal_r$ with the convention
$r_s(x)=\ind\{s(x)>0\}$ inherited from Section~\ref{sec:bg-reject}. By
construction, both $\rho^\theta$ and $c$ are evaluated from the trained
PINN and the solver-cost model alone, with no appeal to $u^\star$, and
Assumption~\ref{ass:linearisation} guarantees that the proxy Bayes rule
$\ind\{\rho^\theta>c\}$ agrees with the true Bayes rule
$\ind\{\eta_h>c\}$ outside a band of width $\mathcal{O}(\delta_{\mathrm{lin}})$
around the level set $\eta_h=c$.

\subsection{The surrogate loss and its pointwise minimiser}
\label{sec:method-surrogate}

The indicator in \eqref{eq:proxy-target} prevents direct gradient-based
minimisation. We replace it with the standard logistic relaxation
\begin{equation}
  \Phi_{\mathrm{abs}}(s;x)
  \;:=\;
  \rho^\theta(x)\,\softplus\!\bigl(-s\bigr)
  \;+\;c(x)\,\softplus\!\bigl(s\bigr),
  \qquad s\in\R,
  \label{eq:surrogate}
\end{equation}
with associated population surrogate risk
$\Ecal_\Phi(s):=\E_X[\Phi_{\mathrm{abs}}(s(X);X)]$. The two terms encode
the two costs the rejector trades: keeping the PINN at $x$ incurs
$\rho^\theta(x)$, and abstaining incurs $c(x)$, each weighted by a smooth
indicator of the corresponding decision. Smoothness of $\softplus$ in
$s$ is what makes \eqref{eq:surrogate} amenable to gradient descent over
$\Hcal_r$.

The first sanity check on \eqref{eq:surrogate} is that its pointwise
minimiser recovers the proxy Bayes rule. This is a direct calibration
fact, analogous to the one familiar from binary logistic regression, but
with a continuous input-dependent ``label'' $\rho^\theta(x)$ and an
input-dependent threshold $c(x)$.

\begin{proposition}[Pointwise minimiser of the surrogate]
\label{prop:pointwise-min}
For $\Pr_X$-a.e.\ $x\in\Omega$ with $\rho^\theta(x),c(x)>0$, the unique
pointwise minimiser of $\Phi_{\mathrm{abs}}(\,\cdot\,;x)$ is
\[
  s^\star(x)\;=\;\log\frac{\rho^\theta(x)}{c(x)},
\]
and the induced rejector $r_{s^\star}(x)=\ind\{s^\star(x)>0\}$ agrees
with $\ind\{\rho^\theta(x)>c(x)\}$, the proxy form of the Bayes rejector
of Lemma~\ref{lem:bayes-rej}. The proof is given in
Appendix~\ref{app:pointwise-min}.
\end{proposition>

Two boundary cases are worth naming. If $\rho^\theta(x)=0$ at some point
the PINN exactly satisfies the local linearised problem and the pointwise
minimiser is $s^\star=-\infty$ (always-PINN); if $c(x)=0$ the solver is
free at that point and $s^\star=+\infty$ (always-solver). Both degenerate
cases are excluded throughout by $\rho^\theta,c>0$ on the active grid;
neither arises in the experiments of Section~\ref{sec:experiments}. In
the regime where Proposition~\ref{prop:pointwise-min} applies, the
surrogate is the unique smooth relaxation of \eqref{eq:proxy-target} whose
pointwise minimiser recovers the proxy threshold rule, and it is the
construction we minimise during training.

\subsection{$\Hcal$-consistency}
\label{sec:method-consistency}

Pointwise calibration is necessary but not sufficient for learning. In
practice $\Ecal_\Phi$ is minimised over a hypothesis class $\Hcal_r$ from
finite samples and finite optimisation budget, and we never reach
$s^\star$ exactly. The relevant guarantee is that small \emph{surrogate}
excess risk $\Ecal_\Phi(s)-\Ecal_\Phi(s^\star)$ implies small
\emph{abstention} excess risk $L(r_s)-L(r_{s^\star})$ within $\Hcal_r$,
in the sense of the $\Hcal$-consistency framework of
\citet{Awasthi_Mao_Mohri_Zhong_2022_multi,
mao2023crossentropylossfunctionstheoretical,Mao_Mohri_Zhong_2023}.
Theorem~\ref{thm:h-consistency} answers in the affirmative with the
square-root rate that is characteristic of strictly convex
surrogates~\citep{bartlett2006convexity}.

\begin{theorem}[$\Hcal$-consistency bound for $\Phi_{\mathrm{abs}}$]
\label{thm:h-consistency}
Assume $\rho_{\max}:=\esssup_{X}\rho^\theta(X)<\infty$ and
$c_{\max}:=\esssup_{X}c(X)<\infty$, and let $\Hcal_r$ contain the
constant scores $\{s\equiv 0\}$. Then for every measurable
$s:\Omega\to\R$,
\begin{equation}
  L(r_s)-L(r_{s^\star})
  \;\leq\;
  \sqrt{\,2\,(\rho_{\max}+c_{\max})\,\bigl(\Ecal_\Phi(s)-\Ecal_\Phi(s^\star)\bigr)\,}.
  \label{eq:h-consistency-main}
\end{equation}
\end{theorem}

The proof, via a Pinsker-type pointwise gap bound on the disagreement set
followed by a Cauchy--Schwarz aggregation, is given in
Appendix~\ref{app:hconsistency}.

The bound \eqref{eq:h-consistency-main} states the operationally relevant
fact: as the rejector class $\Hcal_r$ and the optimisation drive
$\Ecal_\Phi(s)$ toward $\Ecal_\Phi(s^\star)$, the proxy abstention excess
$L(r_s)-L(r_{s^\star})$ vanishes at rate $\sqrt{\,\cdot\,}$. The constant
$2(\rho_{\max}+c_{\max})$ is the natural one: it inherits the heavy-tailed
character of $\rho^\theta$ but is partially compensated because the same
extreme points incur large pointwise surrogate excess. Combined with
Assumption~\ref{ass:linearisation}, the bound controls the \emph{true}
abstention excess up to the linearisation residual
$\delta_{\mathrm{lin}}^\theta$, which is $o(1)$ in the regime where the
PINN is at all useful as a starting point. The hybrid construction of
Section~\ref{sec:hybrid} then turns the rejector's accept/abstain
decisions into a complete PDE field on $\Omega$.
